\documentclass[10pt,a4paper]{amsart}
\usepackage{amsmath,amssymb,amsthm,graphicx}
% Retained minimal preamble for the Stacks Project chapter more-algebra.tex (More on Algebra).
% The chapter file pinned-original-source.tex carries only \input{preamble} on its first line, so no
% preamble text lives inside the chapter itself.  The Stacks Project uses one single shared preamble for
% every chapter; the snapshot stages it identically as sources/stacks-chow/preamble.tex and
% sources/stacks-etale-cohomology/preamble.tex.  The block below is copied mechanically, line for line,
% from sources/stacks-chow/preamble.tex lines 14-18 (the Xy-pic setup) and 154-217 (the native theorem
% declarations together with the mathematical macros).  Nothing was added, removed or redefined; the
% wrapper class and the amsmath/amssymb/amsthm/graphicx base packages are supplied by the assembler.
% Chapter: more-algebra.tex -- More on Algebra.
% For commutative diagrams we use Xy-pic
\usepackage[all]{xy}

% We use 2cell for 2-commutative diagrams.
\xyoption{2cell}

% Theorem environments.
%
\theoremstyle{plain}
\newtheorem{theorem}[subsection]{Theorem}
\newtheorem{proposition}[subsection]{Proposition}
\newtheorem{lemma}[subsection]{Lemma}

\theoremstyle{definition}
\newtheorem{definition}[subsection]{Definition}
\newtheorem{example}[subsection]{Example}
\newtheorem{exercise}[subsection]{Exercise}
\newtheorem{situation}[subsection]{Situation}

\theoremstyle{remark}
\newtheorem{remark}[subsection]{Remark}
\newtheorem{remarks}[subsection]{Remarks}

\numberwithin{equation}{subsection}

% Macros
%
\def\lim{\mathop{\mathrm{lim}}\nolimits}
\def\colim{\mathop{\mathrm{colim}}\nolimits}
\def\Spec{\mathop{\mathrm{Spec}}}
\def\Hom{\mathop{\mathrm{Hom}}\nolimits}
\def\Ext{\mathop{\mathrm{Ext}}\nolimits}
\def\SheafHom{\mathop{\mathcal{H}\!\mathit{om}}\nolimits}
\def\SheafExt{\mathop{\mathcal{E}\!\mathit{xt}}\nolimits}
\def\Sch{\mathit{Sch}}
\def\Mor{\mathop{\mathrm{Mor}}\nolimits}
\def\Ob{\mathop{\mathrm{Ob}}\nolimits}
\def\Sh{\mathop{\mathit{Sh}}\nolimits}
\def\NL{\mathop{N\!L}\nolimits}
\def\CH{\mathop{\mathrm{CH}}\nolimits}
\def\proetale{{pro\text{-}\acute{e}tale}}
\def\etale{{\acute{e}tale}}
\def\QCoh{\mathit{QCoh}}
\def\Ker{\mathop{\mathrm{Ker}}}
\def\Im{\mathop{\mathrm{Im}}}
\def\Coker{\mathop{\mathrm{Coker}}}
\def\Coim{\mathop{\mathrm{Coim}}}

% Boxtimes
%
\DeclareMathSymbol{\boxtimes}{\mathbin}{AMSa}{"02}

%
% Macros for moduli stacks/spaces
%
\def\QCohstack{\mathcal{QC}\!\mathit{oh}}
\def\Cohstack{\mathcal{C}\!\mathit{oh}}
\def\Spacesstack{\mathcal{S}\!\mathit{paces}}
\def\Quotfunctor{\mathrm{Quot}}
\def\Hilbfunctor{\mathrm{Hilb}}
\def\Curvesstack{\mathcal{C}\!\mathit{urves}}
\def\Polarizedstack{\mathcal{P}\!\mathit{olarized}}
\def\Complexesstack{\mathcal{C}\!\mathit{omplexes}}
% \Pic is the operator that assigns to X its picard group, usage \Pic(X)
% \Picardstack_{X/B} denotes the Picard stack of X over B
% \Picardfunctor_{X/B} denotes the Picard functor of X over B
\def\Pic{\mathop{\mathrm{Pic}}\nolimits}
\def\Picardstack{\mathcal{P}\!\mathit{ic}}
\def\Picardfunctor{\mathrm{Pic}}
\def\Deformationcategory{\mathcal{D}\!\mathit{ef}}

% --- appended 2026-10-12 (Chief): the project own definitions of the reference/slogan/history
% display environments, copied mechanically from the shared Stacks preamble lines 8-12
% (\usepackage{verbatim} plus the three \newenvironment lines, which the chapter preamble keeps
% verbatim from the project preamble). Nothing added, removed or redefined: these environments
% carry no mathematical content and the project itself typesets them as comments. Without these
% three lines a counted statement whose own text uses them fails with Environment slogan undefined.
\usepackage{verbatim}
\newenvironment{reference}{\comment}{\endcomment}
\newenvironment{slogan}{\comment}{\endcomment}
\newenvironment{history}{\comment}{\endcomment}
\usepackage{cleveref}
\crefname{originallemma}{original lemma}{original lemmas}
\Crefname{originallemma}{Original lemma}{Original lemmas}
\makeatletter
\newlabel{algebra-lemma-colimit-category-fp-modules}{{05N7}{source}{}{sourcepointer-1}{}}
\newlabel{algebra-lemma-colimit-category-fp-modules@cref}{{[originallemma][0][]05N7}{[1][0][]source}{}{}{}}
\makeatother
\crefname{originallemma}{original lemma}{original lemmas}
\Crefname{originallemma}{Original lemma}{Original lemmas}
\makeatletter
\newlabel{algebra-lemma-flat-finite-presentation-limit-flat}{{02JO}{source}{}{sourcepointer-2}{}}
\newlabel{algebra-lemma-flat-finite-presentation-limit-flat@cref}{{[originallemma][0][]02JO}{[1][0][]source}{}{}{}}
\makeatother
\crefname{originallemma}{original lemma}{original lemmas}
\Crefname{originallemma}{Original lemma}{Original lemmas}
\makeatletter
\newlabel{lemma-structure-relatively-perfect}{{0DHV}{source}{}{sourcepointer-3}{}}
\newlabel{lemma-structure-relatively-perfect@cref}{{[originallemma][0][]0DHV}{[1][0][]source}{}{}{}}
\makeatother
\crefname{originallemma}{original lemma}{original lemmas}
\Crefname{originallemma}{Original lemma}{Original lemmas}
\makeatletter
\newlabel{lemma-compute-RHom-relatively-perfect}{{0E1U}{source}{}{sourcepointer-4}{}}
\newlabel{lemma-compute-RHom-relatively-perfect@cref}{{[originallemma][0][]0E1U}{[1][0][]source}{}{}{}}
\makeatother
\crefname{originallemma}{original lemma}{original lemmas}
\Crefname{originallemma}{Original lemma}{Original lemmas}
\makeatletter
\newlabel{algebra-lemma-directed-colimit-exact}{{00DB}{source}{}{sourcepointer-5}{}}
\newlabel{algebra-lemma-directed-colimit-exact@cref}{{[originallemma][0][]00DB}{[1][0][]source}{}{}{}}
\makeatother
\title{Relative perfectness over a filtered colimit of rings descends to a stage, and Hom groups commute with the colimit}
\author{The Stacks Project Authors}
\date{stacks-project:more-algebra.tex:local-snapshot-20261012; GNU FDL 1.2 or later}
\begin{document}
\maketitle
\setcounter{subsection}{1}
\noindent\textit{Editorial note.} All delivered mathematical text is the author's own, copied from the pinned Stacks Project chapter file more-algebra.tex; no statement or proof step was rewritten, reordered or removed. The counted statement is the labeled lemma environment at source lines 22104-22127 (canonical tag 0DHX) and the counted proof is the proof environment at source lines 22129-22161, which immediately follows it and is its own proof.\par\medskip
\noindent\textit{Editorial note.} The printed text writes that category of R-flat, finitely presented A-modules without the article and begins the last sentence with Thus the result because; both are copied exactly as printed and are not corrected.\par\medskip
\noindent\textit{Editorial note.} The five references are two to results of other chapters on categories of finitely presented modules and one on exactness of filtered colimits, plus two to earlier results of this chapter, among them the Hom complex computation that is itself the subject of a card of this lane. All five are navigation stubs with their roles recorded in the reasons.\par\medskip
\noindent\textit{Editorial note.} The preamble span runs to line 90, the current end of the retained preamble, which includes the appended environment definitions; the counted text does not use them, and the span is taken whole so that the file version is unambiguous.\par\medskip
\noindent\textit{Editorial note.} The source URL is the source-level Stacks Project URL, as required when more than one card is produced from one source; the per-claim tag is carried by claim\_key and by this note. The counted statement and proof cite no bibliography entry, so this card needs no bibliography.\par\medskip
\noindent\textit{Allowed standard prerequisites (editorial).}\begin{itemize}
\item The category of finitely presented modules over a filtered colimit of rings is the colimit of the categories of finitely presented modules over the stages. invoked at proof lines 22129-22132. Reduces the description of finitely presented modules to the stagewise description.
\item The category of flat finitely presented modules is likewise the colimit of the corresponding stagewise categories, and relative perfectness is characterized by a structure statement for such modules. used at proof lines 22133-22140. Produces part (1) by taking a representative of the relatively perfect object at some stage.
\item The Hom complex between a pseudo-coherent complex and a complex of finite tor dimension is a bounded below complex with flat terms whose cohomology computes the Hom groups, and tensor products commute with filtered colimits which are exact. used at proof lines 22142-22160. Computes the Hom groups over the colimit as the colimit of the Hom groups over the stages, which is part (2).
\end{itemize}\medskip
\ifdefined\hypertarget\hypertarget{sourcepointer-1}{}\fi
\noindent\textit{Source reference.} Original lemma 05N7: sources/stacks-chow/tags.txt, source lines 7296--7296. The reference identifies the category of finitely presented modules over the colimit with the colimit of the stagewise categories, invoked at proof lines 22129-22132; the stagewise descent built on it is inside the counted span.\par\medskip
\ifdefined\hypertarget\hypertarget{sourcepointer-2}{}\fi
\noindent\textit{Source reference.} Original lemma 02JO: sources/stacks-chow/tags.txt, source lines 3343--3343. The reference supplies the corresponding description for flat finitely presented modules, used at proof lines 22133-22135; part (1) is derived from it inside the counted span.\par\medskip
\ifdefined\hypertarget\hypertarget{sourcepointer-3}{}\fi
\noindent\textit{Source reference.} Original lemma 0DHV: sources/stacks-more-algebra/pinned-original-source.tex, source lines 21980--21980. The reference is the structure characterization applied at proof lines 22136-22139 to descend the relatively perfect object to a stage; it is a same-chapter pointer and part (1) is completed inside the counted span.\par\medskip
\ifdefined\hypertarget\hypertarget{sourcepointer-4}{}\fi
\noindent\textit{Source reference.} Original lemma 0E1U: sources/stacks-more-algebra/pinned-original-source.tex, source lines 22040--22040. The reference provides the representatives and the Hom complex whose cohomology computes the Hom groups, used at proof lines 22142-22156; the colimit computation built on it is inside the counted span. The lemma itself is the subject of a separate card of this lane with disjoint coordinates.\par\medskip
\ifdefined\hypertarget\hypertarget{sourcepointer-5}{}\fi
\noindent\textit{Source reference.} Original lemma 00DB: sources/stacks-chow/tags.txt, source lines 491--491. The reference gives exactness of filtered colimits, used at proof lines 22158-22160 together with the commutation of tensor products with colimits to identify the group over the colimit; it is the last step and the rest is inside the counted span.\par\medskip
\begin{lemma}
\label{lemma-colimit-relatively-perfect}
Let $R = \colim_{i \in I} R_i$ be a filtered colimit of rings.
Let $0 \in I$ and $R_0 \to A_0$ be a flat ring map of
finite presentation. For $i \geq 0$ set $A_i = R_i \otimes_{R_0} A_0$
and set $A = R \otimes_{R_0} A_0$.
\begin{enumerate}
\item Given an $R$-perfect $K$ in $D(A)$ there exists an $i \in I$
and an $R_i$-perfect $K_i$ in $D(A_i)$ such that
$K \cong K_i \otimes_{A_i}^\mathbf{L} A$ in $D(A)$.
\item Given $K_0, L_0 \in D(A_0)$ with $K_0$ pseudo-coherent
and $L_0$ finite tor dimension over $R_0$, then
we have
$$
\Hom_{D(A)}(K_0 \otimes_{A_0}^\mathbf{L} A, L_0 \otimes_{A_0}^\mathbf{L} A) =
\colim_{i \geq 0}
\Hom_{D(A_i)}(K_0 \otimes_{A_0}^\mathbf{L} A_i,
L_0 \otimes_{A_0}^\mathbf{L} A_i)
$$
\end{enumerate}
In particular, the triangulated category of $R$-perfect complexes over $A$
is the colimit of the triangulated categories of
$R_i$-perfect complexes over $A_i$.
\end{lemma}

\begin{proof}
By Algebra, Lemma \ref{algebra-lemma-colimit-category-fp-modules}
the category of finitely presented $A$-modules is the colimit of
the categories of finitely presented $A_i$-modules.
Given this, Algebra, Lemma
\ref{algebra-lemma-flat-finite-presentation-limit-flat}
tells us that category of $R$-flat, finitely presented $A$-modules
is the colimit of the categories of
$R_i$-flat, finitely presented $A_i$-modules.
Thus the characterization in
Lemma \ref{lemma-structure-relatively-perfect}
proves that (1) is true.

\medskip\noindent
To prove (2) we choose $P_0^\bullet$ representing $K_0$ and
$F_0^\bullet$ representing $L_0$ as in
Lemma \ref{lemma-compute-RHom-relatively-perfect}.
Then $E_0^\bullet = \Hom^\bullet(P_0^\bullet, F_0^\bullet)$
satisfies
$$
H^0(E_0^\bullet \otimes_{R_0} R_i) =
\Hom_{D(A_i)}(K_0 \otimes_{A_0}^\mathbf{L} A_i,
L_0 \otimes_{A_0}^\mathbf{L} A_i)
$$
and
$$
H^0(E_0^\bullet \otimes_{R_0} R) =
\Hom_{D(A)}(K_0 \otimes_{A_0}^\mathbf{L} A, L_0 \otimes_{A_0}^\mathbf{L} A)
$$
by the lemma. Thus the result because tensor product commutes
with colimits and filtered colimits are exact
(Algebra, Lemma \ref{algebra-lemma-directed-colimit-exact}).
\end{proof}

\end{document}
